This project presents a voice-controlled wearable drone wristband system.
Besides conventional RC links, the wristband adds offline speech commands so the operator can command
takeoff, translation, and landing without continuous stick inputs.

Unmanned aerial vehicles (UAVs) controlled through conventional radio transmitters require
the operator to hold a dedicated handheld unit with two joysticks throughout a flight session.
This form factor demands bilateral manual dexterity, prevents the operator from simultaneously
performing other tasks, and creates a barrier for users without prior RC flight experience.
Hands-free or minimal-hardware control interfaces could lower this barrier and enable use
cases in which the operator's hands must remain free, such as first-response inspection or
assistive applications.

Voice command is a well-understood interaction modality that requires no additional hardware
beyond a microphone and a recognition module, imposes no constraint on hand posture, and
maps naturally to a discrete set of high-level flight maneuvers.
This project investigates whether an offline voice-command ground unit can provide sufficient
control fidelity for stable takeoff, directional flight, and landing of a five-inch class quadrotor,
using the ExpressLRS open-source RF link as the wireless bridge.

\subsection{Background and Related Work}

\subsubsection{FPV Quadrotor Platform}

Five-inch class FPV (first-person-view) quadrotors are a well-characterized platform that
balances agility, payload capacity, and component availability.
Standard flight controllers in this class run the Betaflight open-source firmware, which exposes
a CRSF serial interface for RC input and returns sensor telemetry on the same link.
This existing ecosystem provides a stable, well-documented integration target for the ground
control unit.

\subsubsection{ExpressLRS}

ExpressLRS (ELRS) is an open-source 2.4~GHz RC link protocol derived from the LoRa physical
layer \cite{elrs-github}.
It achieves packet update rates from 50~Hz to 1000~Hz with measured round-trip latencies
below 5~ms at 250~Hz, significantly lower than legacy protocols such as Futaba S-FHSS
(55~Hz) or FrSky D16 (150~Hz).
The open firmware stack allows the protocol to run on commodity hardware including the
RadioMaster RP1 and RP2 modules, which use an Espressif ESP8285 microcontroller and a
Semtech SX1280 2.4~GHz transceiver.

\subsubsection{Offline Voice Recognition}

Offline keyword-spotting modules such as the VC02 perform fixed-vocabulary recognition
entirely on-device without network connectivity, achieving recognition latencies typically below
500~ms and operating reliably in environments where cloud-based services are impractical.
The VC02 communicates recognized commands as single-byte UART messages, making
integration with a microcontroller straightforward.

\subsection{System Overview}

\begin{figure}[htb!]
    \centering
    \fbox{\parbox{0.9\linewidth}{\centering
        \vspace{1.8cm}
        \textit{Placeholder block diagram---replace with final vector graphic.}
        \vspace{1.8cm}}}
    \caption{Wearable wristband ground unit, ELRS RF link, and five-inch quadrotor airframe.}
    \label{fig:block-diagram}
\end{figure}

Figure~\ref{fig:block-diagram} shows the complete system.
The ground unit consists of an ESP32 connected to a VC02 voice recognition module and an
SSD1306 OLED display.
The ESP32 encodes operator voice commands as CRSF RC frames and drives the RadioMaster RP2
module via UART.
The RP2, reflashed with ELRS transmitter firmware, broadcasts CRSF frames over the 2.4~GHz ISM band.
On the aircraft, a RadioMaster RX24T receiver forwards decoded RC channels to a SpeedyBee F405~V5
autopilot.
Firmware development began under Betaflight for telemetry experiments and later migrated to
ArduPilot/ArduCopter so altitude hold could run on-board (Section~\ref{sec:flight-controller}).
Four 1860~KV brushless motors are driven by an OX32 55~A four-in-one ESC\@.
Telemetry (including barometric altitude when scheduling permits) returns over ELRS for OLED cues on
the wrist.

\subsection{High-Level Requirements}

Table~\ref{tab:hlr} lists the top-level performance requirements for the system.

\begin{table}[htb!]
    \centering
    \caption{High-level system requirements.}
    \label{tab:hlr}
    \begin{tabularx}{\linewidth}{l X l}
        \hline
        \textbf{Requirement} & \textbf{Description} & \textbf{Target} \\
        \hline
        R1 & Voice command recognition latency from utterance end to VC02 UART byte &
            $<$1~s \\
        R2 & End-to-end command latency from VC02 byte to CRSF frame departure at RP2 &
            $<$40~ms \\
        R3 & RF link update rate during flight & $\geq$50~Hz \\
        R4 & Stable hover after takeoff command without pilot intervention &
            $\geq$5~s \\
        R5 & Directional movement response to forward/backward command & Observable
            displacement \\
        R6 & Successful autonomous landing (disarm without crash) & 100\% of attempts \\
        R7 & Barometric altitude displayed on OLED during flight & Continuous \\
        \hline
    \end{tabularx}
\end{table}
